% Bagian: first-order-logic
% Bab: model-theory
% Subbagian: partial-iso

\documentclass[../../../include/open-logic-section]{subfiles}

\begin{document}

\olfileid[id]{mod}{bas}{dlo}
\section{Urutan Linear Rapat}

\begin{defn}
  Suatu \emph{urutan linear rapat tanpa titik ujung} ialah
  !!a{structure} $\Struct{M}$ bagi !!{language} yang memuat satu
  !!{predicate}~$<$ beraritas~2 dan memenuhi kalimat-kalimat berikut:
  \begin{enumerate}
  \item $\lforall[x][\lnot x < x]$;
  \item $\lforall[x][\lforall[y][\lforall[z][(x < y \lif (y < z \lif x
    <z ))]]]$;
  \item $\lforall[x][\lforall[y][(x< y \lor \eq[x][y] \lor y < x)]]$;
  \item $\lforall[x][\lexists[y][x < y]]$;
  \item $\lforall[x][\lexists[y][y < x]]$;
  \item $\lforall[x][\lforall[y][(x < y \lif \lexists[z][(x < z \land
        z < y))]]]$.
 \end{enumerate}
\end{defn}

\begin{thm}\ollabel{thm:cantorQ}
  Sembarang dua urutan linear rapat tanpa titik ujung yang !!{enumerable}
  saling isomorfik.
\end{thm}

\begin{proof}
  Misalkan $\Struct{M_1}$ dan $\Struct{M_2}$ merupakan urutan linear rapat
  tanpa titik ujung yang !!{enumerable}, dengan ${<_1} = \Assign{<}{M_1}$
  dan ${<_2} = \Assign{<}{M_2}$, serta misalkan $\PIso{I}$ merupakan
  himpunan semua isomorfisme parsial di antara keduanya. $\PIso{I}$ tidak
  kosong karena sekurang-kurangnya $\emptyset \in \PIso{I}$. Kita tunjukkan
  bahwa $\PIso{I}$ memenuhi sifat bolak-balik. Dengan demikian,
  $\Struct{M_1} \iso[p] \Struct{M_2}$, dan teorema ini mengikuti dari
  \olref[pis]{thm:p-isom1}.

  Untuk menunjukkan bahwa $\PIso{I}$ memenuhi sifat Maju, misalkan
  $p \in \PIso{I}$. Apabila fungsi parsial ini tidak kosong, tulis
  $p(a_i) = b_i$ bagi $i = 1$, \dots,~$n$. Tanpa mengurangi keumuman,
  andaikan $a_1 <_1 a_2 <_1 \cdots <_1 a_n$.
  Diberikan $a \in \Domain{M_1}$, kita memperoleh perluasan yang diinginkan
  sebagai berikut. Jika unsur yang diberikan sudah berada dalam domain
  fungsi parsial itu, fungsi tersebut sendiri merupakan perluasan yang
  diinginkan. Jika fungsi parsial itu kosong, pasangan tunggal apa pun dengan
  unsur yang diberikan sebagai komponen pertama dan suatu anggota struktur
  kedua sebagai komponen kedua memberikan perluasan yang diinginkan. Dalam
  semua kasus lainnya, unsur yang diberikan tidak berada dalam domain fungsi
  parsial itu dan kita menemukan $b \in \Domain{M_2}$ sebagai berikut:
  \begin{enumerate}
  \item jika $a <_1 a_1$, ambil $b \in \Domain{M_2}$ sedemikian sehingga
    $b <_2 b_1$;
  \item jika $a_n <_1 a$, ambil $b \in \Domain{M_2}$ sedemikian sehingga
    $b_n <_2 b$;
  \item jika $a_i <_1 a <_1 a_{i+1}$ bagi suatu $i$, ambil
    $b \in \Domain{M_2}$ sedemikian sehingga $b_i <_2 b <_2 b_{i+1}$.
  \end{enumerate}
  Selalu mungkin menemukan $b$ dengan sifat yang diinginkan karena
  $\Struct{M_2}$ merupakan urutan linear rapat tanpa titik ujung. Dalam
  kasus-kasus terakhir ini, tetapkan $q = p \cup \{ \langle a, b \rangle
  \}$; maka $q \in \PIso{I}$ merupakan perluasan $p$ yang diinginkan. Ini
  menetapkan sifat Maju. Sifat Mundur dibuktikan dengan cara serupa. Jadi,
  $\Struct{M_1} \iso[p] \Struct{M_2}$; menurut
  \olref[pis]{thm:p-isom1}, $\Struct{M_1} \iso \Struct{M_2}$.
\end{proof}

\begin{prob}
  Lengkapi bukti \olref[mod][bas][dlo]{thm:cantorQ} dengan memverifikasi
  bahwa $\PIso{I}$ memenuhi sifat Mundur.
\end{prob}

\begin{rem}
  Misalkan $\Struct{S}$ sembarang urutan linear rapat tanpa titik ujung yang
  !!{enumerable}. Maka (menurut \olref{thm:cantorQ}) $\Struct{S} \iso
  \Struct{Q}$, dengan $\Struct{Q} = (\Rat, <)$ sebagai urutan linear rapat
  !!{enumerable} yang domainnya merupakan himpunan bilangan rasional~$\Rat$.
  Sekarang tinjau kembali !!{structure}~$\Struct{R} = (\Real, <)$ dari
  \olref[thm]{remark:R}. Kita telah melihat bahwa terdapat
  !!a{structure}~$\Struct{S}$ yang !!{enumerable} sedemikian sehingga
  $\Struct{R} \elemequiv \Struct{S}$. Namun, $\Struct{S}$ merupakan
  urutan linear rapat tanpa titik ujung yang !!{enumerable}, sehingga struktur itu
  isomorfik (dan karenanya ekuivalen secara elementer) dengan
  !!{structure}~$\Struct{Q}$. Menurut transitivitas ekuivalensi elementer,
  $\Struct{R} \elemequiv \Struct{Q}$. (Kita juga dapat menunjukkan ini
  secara langsung dengan menetapkan $\Struct{R} \iso[p] \Struct{Q}$ melalui
  argumen bolak-balik yang sama.)
\end{rem}
\end{document}
